\section{Conclusion}
\label{sec:conclusion}
% Le & Trung's text
%This work investigated the so-called selfish scheduling phenomenon, which degrades the performance of MPQUIC schedulers in the multi-client context. 
%for wireless networks with multiple clients communicating with a sever simultaneously. 
%To this end, we proposed MuLeS, a novel MPQUIC scheduler based on Q-learning that allows for scheduling packets for multiple clients simultaneously.
%By taking states of all sessions into account, MuLeS can offer a scheduling decision fairly for all clients, thereby alleviating the selfish scheduling phenomenon. Consequently, MuLeS can shorten the mean download time of all clients by up to 16\% compared to other schedulers.
%However, the performance of MuLeS depends on the values of some hyperparameters, such as $\beta$. Our future work will be devoted to automating the selection of optimal hyperparameters, making our algorithm more robust.


%Kien's text 2023 Sep. 3
This work investigated the MPQUIC scheduling issue with the consideration of multiple clients. We then proposed MuLeS, a novel MPQUIC scheduler based on Q-learning that allows for simultaneously scheduling packets for several clients. MuLes considers the states of all concurrent sessions and offers a scheduling decision fairly for all clients. As a result, MuLeS alleviates the rush scheduling phenomenon that happened with the existing schedulers. We have implemented and evaluated MuLes in comparison to state-of-the-art MPQUIC schedulers. The results show that  MuLeS can shorten the mean download time of all clients by up to 16\% compared to the others. However, the performance of MuLeS depends on the values of some hyperparameters, such as $\beta$. Our future work will be devoted to automating the selection of optimal hyperparameters, making our algorithm more robust.
%\balance